\chapter{Alternativas consideradas}

Se compararon seis familias de robots o vehiculos moviles. La comparacion se centra en la idoneidad para una nave aeronautica interior con operarios, utiles grandes, estaciones cambiantes y necesidad de trazabilidad.

\begin{longtable}{P{0.20\textwidth}P{0.30\textwidth}P{0.18\textwidth}P{0.22\textwidth}}
\toprule
\textbf{Alternativa} & \textbf{Ventajas} & \textbf{Limitaciones} & \textbf{Decision}\\
\midrule
\endhead
AGV & Robusto para rutas repetitivas, tecnologia madura, coste predecible. & Menor flexibilidad, dependencia de rutas fijas o infraestructura, dificil adaptacion a cambios de layout. & Descartado como opcion principal; valido solo para corredores muy estables.\\
AMR diferencial & Navegacion flexible, buena disponibilidad comercial, integracion con flota. & Maniobra peor en espacios estrechos que una base omnidireccional. & Alternativa viable, pero no optima para estaciones congestionadas.\\
AMR omnidireccional & Movimiento lateral, alta maniobrabilidad, buena convivencia con estaciones y utiles grandes. & Mayor coste y complejidad mecanica que una base diferencial sencilla. & \textbf{Seleccionado}.\\
Robot con orugas & Buen contacto en terrenos irregulares y capacidad de superar obstaculos. & Excesivo para suelo industrial interior, peor higiene, posible daño al pavimento, maniobra poco fina. & Descartado.\\
Dron & Rapido para inspeccion visual en altura. & Riesgo de seguridad, ruido, autonomia limitada, restricciones en interior industrial y carga practicamente nula. & Descartado para logistica y FOD en estaciones.\\
Robot con patas & Alta movilidad teorica en entornos complejos. & Coste alto, complejidad, riesgo de aceptacion, innecesario en nave con pavimento regular. & Descartado.\\
\bottomrule
\end{longtable}

\section{Robot seleccionado}

La solucion adopta una \textbf{plataforma AMR omnidireccional}, representable en CoppeliaSim mediante KUKA YouBot, KUKA Omnirob o una base holonomica equivalente. Esta eleccion es coherente con la asignatura porque permite trabajar con modelos de robot movil, sensores de navegacion, planificacion de trayectorias y validacion en simulador antes de proponer un despliegue fisico.

La base omnidireccional aporta tres ventajas concretas: puede alinearse con estaciones sin giros amplios, puede corregir lateralmente su posicion al aproximarse a utiles o mesas, y reduce maniobras innecesarias en zonas compartidas con operarios. En un entorno aeronautico interior, estas ventajas pesan mas que la simplicidad de un AGV de ruta fija.
